\par\Needspace{6\baselineskip}\medskip
\hypertarget{native-207-article153}{}\pdfbookmark[2]{153.}{article-153}
\begin{center}153.\end{center}
In the \textit{second hypothesis} we shall assign to $P$, $Q$ the very values that we found for $P'$, $Q'$ in the first. We shall therefore set
\[
\begin{aligned}
x&=\log P=0{,}079\,0164\\
y&=\log Q=8{,}5475981.
\end{aligned}
\]

Since the computation is to be carried out here in exactly the same way as in the first hypothesis, it will suffice to give its principal results:
\begin{center}
\small
\setlength{\tabcolsep}{6pt}
\begin{tabular}{@{}r|r@{}}
$\omega=13^0\,15'\,38''{,}13$ & $\zeta''=210^0\,8'\,24''{,}98$\\[.5em]
$\omega+\sigma=13\;\;38\;\;51{,}25$ & $\log r=0{,}330\,7676$\\[.5em]
$\log Qc\sin\omega=0{,}598\,9389$ & $\log r''=0{,}322\,2280$\\[.5em]
$z=14^0\,33'\,19''{,}00$ & $\tfrac12(u''+u)=205^0\,22'\,15''{,}58$\\[.5em]
$\log r'=0{,}325\,9918$ & $\tfrac12(u''-u)=-3\;\;14\;\;4{,}79$\\[.7em]
$\log\dfrac{n'r'}{n}=0{,}667\,5193$ & $2f'=7\;\;34\;\;53{,}32$\\[.7em]
$\log\dfrac{n'r'}{n''}=0{,}588\,5029$ & $2f=3\;\;29\;\;0{,}18$\\[.5em]
$\zeta=203^0\,16'\,38''{,}16$ & $2f''=4\;\;5\;\;53{,}12$
\end{tabular}
\end{center}

It would hardly be worth recomputing the reductions of the times for aberration, since they differ by scarcely $2^{\mathrm{s}}$ from those we obtained in the first hypothesis.

Further computations give $\log\eta=0{,}000\,2270$, $\log\eta''=0{,}000\,3173$, from which we derive
\SourceLeaf{208}
\[
\begin{aligned}
\log P'&=0{,}079\,0167,&\qquad X&=+0{,}000\,0003\\
\log Q'&=8{,}5476110,&\qquad Y&=+0{,}000\,0129.
\end{aligned}
\]
This shows how much more accurate the second hypothesis is than the first.
